\subsection{Projector-to-hardware falsifier: from graph idempotents to timing ports}
\label{sec:projector-hardware-falsifier}

The clock--matter resonance is now not only a spectral coincidence but an explicit
projector-to-hardware test chain.  The rank-$120$ resonance and rank-$24$
companion blocks are separated by polynomial graph projectors.  For the Heawood
flag-clock Laplacian $L_c$ and matter Laplacian $L_m$, the minimal monomial
compiler uses
\[
P_{c,6}=\frac{1}{42}L_c^3-\frac17L_c^2+\frac16L_c,
\qquad
P_{c,0}=I-\frac{43}{42}L_c+\frac27L_c^2-\frac{1}{42}L_c^3,
\]
and
\[
P_{m,24}=\frac{5}{24}L_m-\frac{1}{144}L_m^2,
\qquad
P_{m,30}=-\frac{2}{15}L_m+\frac{1}{180}L_m^2.
\]
Thus the minimal graph hardware only needs clock walk powers through $L_c^3$ and
matter walk powers through $L_m^2$.  The two analyzer ports are
\[
P_{\rm res}=P_{c,6}\otimes P_{m,24},\qquad \mathrm{rank}=120,
\]
and
\[
P_{\rm comp}=P_{c,0}\otimes P_{m,30},\qquad \mathrm{rank}=24.
\]
Their raw minimal-degree LCU masses are
\[
\|P_{\rm res}\|_1=\frac{31}{432},\qquad
\|P_{\rm comp}\|_1=\frac{35}{108},\qquad
\|P_{\rm res}\|_1+\|P_{\rm comp}\|_1=\frac{19}{48},
\]
with maximum walk depth $5$ inside the $2048$-bin envelope.

A first component-named optical budget assigns placeholder survival values to
switches, delays, phase shifters, analyzers, detectors, and dark-count background.
These are explicitly not measured hardware constants.  Under the default
placeholder budget the two selector ports have shot-noise SNRs
\[
S_{\rm res}=36.666463881272,
\qquad
S_{\rm comp}=36.791565396791,
\]
above the five-sigma rule $S\ge5$.  A component-range sweep over $960$ cases gives
$721$ passes and $239$ failures; all cases with parity/sign-flip probability
$p_{\rm flip}\le0.20$ pass, whereas all $p_{\rm flip}=0.45$ cases fail.  Thus the
separator is primarily parity/sign-flip limited, not uniformly loss limited.

The LCU optimizer now has three layers.  Raw algebraic coefficients favor high
monomial degree; the raw tested frontier prefers $(d_c,d_m)=(9,8)$ with
\[
\|c\|_{1,\rm raw}=2.0822330410596202\times10^{-10}.
\]
But block-encoded hardware normally supplies $H=L/\Lambda$, not $L$ itself.  The
normalization rule
\[
\sum_i c_iL^i=\sum_i c_i\Lambda^iH^i
\]
changes the hardware cost.  With
\[
\Lambda_c=6,
\qquad
\Lambda_m=30,
\]
the raw $(9,8)$ point becomes
\[
\|c\|_{1,\rm block}=289713.4069956163,
\]
while the best tested monomial block-encoded point is instead
\[
(d_c,d_m)=(4,2),
\qquad
\|c\|_{1,\rm block}=334.6461794019932.
\]

The normalized Chebyshev/QSVT route has an important parity caveat.  In the
centered signal model
\[
x=2(L/\Lambda)-1,
\]
single-sequence QSVT polynomials obey
\[
p(-x)=(-1)^d p(x).
\]
Therefore endpoint selectors such as $P_{c,6}$, $P_{c,0}$, and $P_{m,30}$ cannot
be implemented by a single parity-constrained sequence: they distinguish
$x=-1$ from $x=1$.  They instead require a two-sequence even/odd LCU or
ancilla-selected QSVT route.

For the clock endpoint selectors, with clock spectral points
\[
\{-1,-\sqrt2/3,\sqrt2/3,1\},
\]
use
\[
e_c(x)=\frac{9}{14}x^2-\frac17
=\frac{5}{28}T_0(x)+\frac{9}{28}T_2(x),
\]
and
\[
o_c(x)=\frac{9}{14}x^3-\frac17x
=\frac{19}{56}T_1(x)+\frac{9}{56}T_3(x).
\]
Then
\[
P_{c,6}=e_c+o_c,
\qquad
P_{c,0}=e_c-o_c,
\]
with two-sequence mass $1$ for either clock endpoint selector.  For the matter-30
endpoint selector, use
\[
e_{30}(x)=\frac54x^2-\frac34,
\qquad
o_{30}(x)=x/2,
\]
so $P_{m,30}=e_{30}+o_{30}$ has two-sequence mass $1.25$.

The matter-24 projector is the positive single-sequence case.  It is certified by
the even quartic
\[
P_{m,24}(x)=-\frac{625}{256}x^4+\frac{225}{128}x^2+\frac{175}{256},
\]
where $x=2(L_m/30)-1$.  This polynomial satisfies
\[
P_{m,24}(-1)=0,
\qquad
P_{m,24}(3/5)=1,
\qquad
P_{m,24}(1)=0,
\]
and has analytic certificate
\[
\|P_{m,24}\|_{\infty,[-1,1]}=1.
\]
The resulting two-port logical QSVT/LCU mass is
\[
\boxed{2.5439453125},
\]
namely $1.2939453125$ for the resonance port and $1.25$ for the companion port.
The Chebyshev-term phase schedules are exact at the term level:
\[
T_0:[],\quad T_1:[0],\quad T_2:[0,0],\quad T_3:[0,0,0],\quad T_4:[0,0,0,0].
\]
The whole-polynomial implementation is therefore a certified LCU of QSP terms;
collapsed single-sequence phase lists are not claimed.

The parity-routed hardware resource table has max tensor depth $7$.  The
resonance port has mass $1.2939453125$, $12$ term products, and placeholder SNR
$37.034240801677326$.  The companion port has mass $1.25$, $12$ term products,
and placeholder SNR $37.33164132928002$.  Abstract LCU success accounting gives
\[
p_{\rm res}=0.597413338718975,
\qquad
p_{\rm comp}=0.64.
\]
The corresponding unamplified effective SNRs are
\[
28.625716516580017,
\qquad
29.865313063424016.
\]
A monolithic two-port LCU with mass $2.5439453125$ has success probability
\[
0.1544685004281349
\]
and effective SNR
\[
14.557904134592752.
\]
Thus separate-port postselection is already high-success under placeholder
assumptions, while a monolithic two-port selection should use extra shots or
amplification in an actual run card.

Phase precision reinforces the same conclusion.  With the first-order bound
\[
\epsilon_{\rm proj}\le \kappa\sigma_\phi,
\]
the block-encoded $(4,2)$ point tolerates
\[
\sigma_{\phi,\max}=1.71079969939437\times10^{-3}
\]
for a $10^{-2}$ projector-error budget, whereas the high-degree $(8,8)$ and
$(9,8)$ candidates require
\[
4.3872988575743546\times10^{-7}
\quad\text{and}\quad
1.0112501159704773\times10^{-7},
\]
respectively.  Thus the near-term optical target is shallow normalized or
parity-routed QSVT compilation, not raw high-degree monomial interpolation.

The clock--Levi bridge also has a resolved boundary.  A chosen spanning-tree/cycle
basis gives an injective coordinate bridge $H_1(H)\to H_1(L_{W33})$ of rank $8$,
but the bridge is gauge-fixed rather than canonical.  Varying the Levi root over
all $80$ vertices gives $38$ distinct support signatures.  Root-gauge twirling
spreads the eight-cycle selector over all $160$ Levi incidence edges and all $80$
Levi vertices.  The projective symplectic action generated by transvections has
order $25920$ (the $\pm I$ center of $\mathrm{Sp}(4,3)$ acts trivially
projectively), and the orbit of a single Levi incidence edge has size $160$.
Hence the automorphism twirl is the uniform incidence-edge idempotent, not a
hidden sparse eight-cycle embedding.

The uniform incidence-edge idempotent has no protected Levi cycle content.  For
the W33 Levi chain complex,
\[
|V|=80,
\qquad
|E|=160,
\qquad
\operatorname{rank}D=79,
\qquad
\beta_1=81.
\]
But if $E_{\rm unif}=\mathbf1_E\mathbf1_E^T/160$, then
\[
\|P_{H_1}\mathbf1_E\|=2.3633435577592544\times10^{-14},
\qquad
\frac{\mathbf1_E^TP_{H_1}\mathbf1_E}{\mathbf1_E^T\mathbf1_E}
=-6.817896941457846\times10^{-16}.
\]
Thus full support twirling kills the homological bridge.

For an oriented bridge subspace, however, the automorphism-twirl result is the
opposite and exactly what Schur's lemma predicts.  Let $P_B$ be the projector onto
a fixed oriented 8-cycle bridge subspace inside $H_1$.  The automorphism average
is
\[
\frac1{|G|}\sum_{g\in G}\rho(g)P_B\rho(g)^{-1}
=\frac{8}{81}P_{H_1}.
\]
The exact character check removes the previous irreducibility caveat.  For each
group element,
\[
\chi_{H_1}(g)=f_E(g)-f_V(g)+1,
\]
and the value distribution
\[
-3^{(810)},\ -1^{(3240)},\ 0^{(16640)},\ 1^{(5184)},\ 9^{(45)},\ 81^{(1)}
\]
has square sum $25920$.  Hence
\[
\langle\chi_{H_1},\chi_{H_1}\rangle=1,
\]
so the $81$-dimensional Levi $H_1$ character is irreducible over $\mathbb C$ for
the generated projective symplectic action.  Full symmetry averaging erases the
chosen bridge basis but preserves its oriented $H_1$ trace as isotropic density.
This reconciles the support-twirl and subspace-twirl stories.

The falsifier chain is therefore precise.  First build the graph projectors as
normalized walk-power or parity-routed QSVT/LCU ports.  Then run the $30$-strobe
and clock-sector separator from \S\ref{sec:heawood-clock-homology}.  The expected
timing readout is still
\[
P_{c,6}\otimes P_{m,24}\mapsto -1,
\qquad
P_{c,0}\otimes P_{m,30}\mapsto +1
\]
under the partial clock evolution at $\tau=\pi/6$.  A failure of this split,
after calibrated loss, sign-flip, block-encoding normalization, QSVT parity
routing, ancilla success accounting, and phase-precision budgets are applied,
falsifies the projector version of the clock--matter resonance.  (Witnesses:
\texttt{analysis/bt1661\_projector\_hardware\_compiler.py},
\texttt{analysis/bt1667\_calibrated\_loss\_sweep.py},
\texttt{analysis/bt1673\_block\_encoding\_normalization\_audit.py},
\texttt{analysis/bt1679\_parity\_valid\_qsvt\_phase\_compiler.py},
\texttt{analysis/bt1680\_analytic\_sup\_norm\_certificate.py},
\texttt{analysis/bt1685\_qsp\_phase\_angle\_synthesis.py},
\texttt{analysis/bt1688\_exact\_h1\_character\_irreducibility.py}, and
\texttt{analysis/bt1689\_ancilla\_success\_accounting.py}.)
